\documentclass[final]{cvpr}

% Packages
\usepackage{graphicx}
\usepackage{amsmath}
\usepackage{amssymb}
\usepackage{booktabs}
\usepackage{multirow}
\usepackage{algorithm}
\usepackage{algorithmic}
\usepackage{hyperref}
\usepackage{xcolor}
\usepackage{subcaption}
\usepackage[utf8]{inputenc}
\usepackage[T1]{fontenc}

\definecolor{cvprblue}{rgb}{0.21,0.49,0.74}
\hypersetup{colorlinks,linkcolor={cvprblue},citecolor={cvprblue},urlcolor={cvprblue}}

\def\cvprPaperID{****}
\def\confName{CVPR}
\def\confYear{2026}

\title{Person Re-Identification: From Classical Descriptors to Deep Metric Learning with Application to Photo Album Clustering}

\author{
Maxence Rossignol\\
Daphné Maschas\\
Adham Noureldin\\
{\tt\small Final Project Report}\\
{\tt\small Visual Recognition \& Deep Learning}
}

\begin{document}
\maketitle

% ============================================================
% ABSTRACT
% ============================================================
\begin{abstract}
Person Re-Identification (Re-ID) aims to match images of the same individual across non-overlapping camera views and is critical for video surveillance and automated photo organization.
In this work, we present a comprehensive study of person Re-ID on the Market-1501 benchmark, spanning classical computer vision descriptors, deep metric learning architectures, and unsupervised clustering for album indexing.
We first analyze the impact of image preprocessing, including color space representations, spatial filtering, data augmentation strategies, and input resolution on downstream Re-ID accuracy.
We then systematically compare classical feature extractors---HOG, color histograms, and SIFT---against a ResNet-50 backbone equipped with a Batch Normalization Neck (BN-Neck) trained with a joint cross-entropy and triplet loss.
An ablation study on backbone freezing strategies reveals that full fine-tuning is essential, achieving 85.54\% Rank-1 accuracy and 68.23\% mAP after a single training epoch.
We further demonstrate that k-reciprocal re-ranking provides a +9.66\% mAP improvement by exploiting neighborhood consistency.
Finally, we apply the learned embeddings to unsupervised identity clustering using DBSCAN, achieving 88.95\% average album purity, demonstrating practical applicability for automatic photo album organization.
Our experiments comprehensively bridge classical visual recognition with modern deep learning for the Re-ID task.
\end{abstract}

% ============================================================
% 1. INTRODUCTION
% ============================================================
\section{Introduction}

Person Re-Identification (Re-ID) is the task of recognizing an individual across different camera views, times, and conditions.
It is a fundamental problem in computer vision with immediate applications in video surveillance~\cite{zheng2016person}, forensic investigation, and smart city infrastructure.
Beyond security, Re-ID has emerging applications in consumer technology, such as automatically grouping photos of the same person in a personal photo library---analogous to face clustering but operating on full-body appearance.

The challenge of person Re-ID lies in the significant intra-class variability and inter-class similarity present in real-world imagery.
The same individual may appear drastically different across cameras due to changes in illumination, viewpoint, pose, and partial occlusions.
Conversely, different individuals wearing similar clothing can appear nearly identical.
These challenges distinguish Re-ID from classical image retrieval and demand representations that are both discriminative and invariant to nuisance factors.

Traditional approaches relied on handcrafted features---Histogram of Oriented Gradients (HOG)~\cite{dalal2005hog}, Scale-Invariant Feature Transform (SIFT)~\cite{lowe2004sift}, and color histograms---combined with distance metrics.
While effective for instance-level matching, these methods struggle with the viewpoint and illumination variations inherent to multi-camera networks.
The advent of deep learning has revolutionized Re-ID: convolutional neural networks (CNNs) learn rich, transferable representations that dramatically outperform handcrafted features~\cite{he2016deep,luo2019bag}.

In this project, we conduct a comprehensive study of person Re-ID that bridges classical visual recognition and modern deep learning.
Our contributions are as follows:
\begin{itemize}
    \item A systematic analysis of image preprocessing choices---color spaces, filtering, augmentation, and resolution---and their impact on Re-ID performance.
    \item A quantitative comparison between classical descriptors (HOG, color histograms, SIFT) and deep embeddings from a ResNet-50 backbone with BN-Neck.
    \item An ablation study on transfer learning strategies (feature extraction, partial fine-tuning, full fine-tuning).
    \item An evaluation of k-reciprocal re-ranking as post-processing for improving retrieval quality.
    \item A practical application of learned embeddings for unsupervised photo album clustering using DBSCAN.
\end{itemize}

% ============================================================
% 2. PROBLEM DEFINITION
% ============================================================
\section{Problem Definition}

\subsection{Notation and Definitions}
Let $\mathcal{D} = \{(x_i, y_i, c_i)\}_{i=1}^{N}$ denote a person Re-ID dataset where $x_i \in \mathbb{R}^{H \times W \times 3}$ is an RGB image, $y_i \in \{1, \ldots, K\}$ is the identity label, and $c_i \in \{1, \ldots, C\}$ is the camera index.
The dataset is partitioned into a training set $\mathcal{D}_{\text{train}}$, a query set $\mathcal{Q}$, and a gallery set $\mathcal{G}$ with identity-disjoint splits in the standard protocol.

\subsection{Objective}
Given a query image $x_q$ with identity $y_q$, the goal is to retrieve all images in the gallery $\mathcal{G}$ depicting the same identity.
Formally, we seek a feature extractor $f_\theta : \mathbb{R}^{H \times W \times 3} \rightarrow \mathbb{R}^d$ parameterized by $\theta$ such that:
\begin{equation}
    d(f_\theta(x_q), f_\theta(x_g)) < d(f_\theta(x_q), f_\theta(x_n))
\end{equation}
for all gallery images $x_g$ with $y_g = y_q$ and $x_n$ with $y_n \neq y_q$, where $d(\cdot, \cdot)$ is a distance metric (Euclidean or cosine).

\subsection{Evaluation Protocol}
Following the standard Market-1501 protocol~\cite{zheng2015market}, gallery images sharing both the same identity \emph{and} the same camera as the query are excluded from evaluation (trivially easy matches).
Performance is measured by:
\begin{itemize}
    \item \textbf{Rank-$k$ accuracy}: the probability that at least one correct match appears in the top-$k$ retrieved results.
    \item \textbf{Mean Average Precision (mAP)}: the mean of per-query average precision, measuring overall retrieval quality.
    \item \textbf{Cumulative Matching Characteristic (CMC)}: the cumulative curve of Rank-$k$ for $k = 1, \ldots, 50$.
\end{itemize}

\subsection{Constraints}
The system must handle:
(i)~cross-camera illumination and color shifts,
(ii)~pose and viewpoint variations,
(iii)~partial occlusions from carried objects or scene elements,
(iv)~background clutter within bounding boxes, and
(v)~resolution variability from differing camera-to-subject distances.

% ============================================================
% 3. RELATED WORK
% ============================================================
\section{Related Work}

\textbf{Handcrafted Features for Re-ID.}
Early person Re-ID systems relied on handcrafted descriptors.
Farenzena~\emph{et al.}~\cite{farenzena2010person} proposed the Symmetry-Driven Accumulation of Local Features (SDALF) combining HSV color histograms with texture features.
Gray and Tao~\cite{gray2008viewpoint} used an ensemble of color and texture features with boosting for cross-view matching.
While these approaches established the field, they lack the invariance needed for large-scale multi-camera settings.

\textbf{Deep Metric Learning.}
The transition to deep learning for Re-ID was catalyzed by works such as~\cite{yi2014deep}, which first applied Siamese networks to learn discriminative embeddings.
Hermans~\emph{et al.}~\cite{hermans2017defense} demonstrated that batch-hard triplet mining significantly improves metric learning for Re-ID.
The triplet loss encourages the network to produce embeddings where images of the same person are closer than those of different persons by at least a margin $\alpha$.

\textbf{BN-Neck Architecture.}
Luo~\emph{et al.}~\cite{luo2019bag} introduced a strong baseline combining a ResNet-50 backbone with a Batch Normalization Neck (BN-Neck).
The key insight is to apply the classification (cross-entropy) loss on features \emph{after} batch normalization, while applying the triplet loss on features \emph{before} normalization.
This decoupling allows the two losses to optimize complementary objectives.
Our implementation follows this architecture closely.

\textbf{Re-ranking Methods.}
Zhong~\emph{et al.}~\cite{zhong2017reranking} proposed k-reciprocal re-ranking, a post-processing method that refines the initial ranking by exploiting the similarity structure among retrieved results.
The method constructs a Jaccard distance from k-reciprocal nearest neighbor sets and combines it with the original Euclidean distance.
This approach is architecture-agnostic and consistently improves mAP across different feature extractors, which we confirm in our experiments.

\textbf{Transformer-based Re-ID.}
Recent work has explored Vision Transformers (ViT) for Re-ID~\cite{he2021transreid}, leveraging self-attention to capture global dependencies.
While ViTs show promising results, they typically require larger training budgets.
Our study focuses on the CNN-based ResNet-50 baseline as the primary architecture.

\textbf{Unsupervised Clustering for Re-ID.}
When identity labels are unavailable, clustering-based methods group person images based on embedding similarity.
Lin~\emph{et al.}~\cite{lin2019bottom} proposed a bottom-up clustering framework for unsupervised Re-ID.
DBSCAN~\cite{ester1996dbscan} is particularly suitable for this task as it does not require specifying the number of clusters \emph{a priori} and can handle noise points.
We leverage this approach for our photo album indexing application.

% ============================================================
% 4. METHODOLOGY
% ============================================================
\section{Methodology}

\subsection{Data and Preprocessing}

\textbf{Market-1501 Dataset.}
We use the Market-1501 benchmark~\cite{zheng2015market}, comprising 32,668 bounding box images of 1,501 identities captured by 6 cameras.
The training set contains 12,936 images of 751 identities.
The query and gallery sets contain 3,368 and 13,115 images respectively, covering 750 identities.
All images are person bounding boxes detected by a DPM detector~\cite{felzenszwalb2010dpm}.

\textbf{Color Space Analysis.}
We study three color representations: RGB, HSV, and CIE-LAB.
The HSV space decouples chromatic information (Hue, Saturation) from luminance (Value), providing partial illumination invariance.
The LAB space is designed to be perceptually uniform, with the $L$ channel encoding lightness and the $a$, $b$ channels encoding color opponency.
Our analysis reveals that the Hue channel produces stable identity signatures---the dominant clothing color appears as sharp peaks---while RGB channels exhibit high inter-channel correlation and sensitivity to lighting changes.

\textbf{Spatial Filtering.}
We evaluate Gaussian, median, and bilateral filters with kernel size $k=5$.
The bilateral filter achieves the best trade-off between noise reduction and edge preservation: it smooths uniform areas while maintaining sharp object boundaries, which is critical for preserving identity-discriminative silhouette and texture features.

\textbf{Data Augmentation.}
We implement three augmentation strategies:
(i)~\emph{Standard}: random horizontal flip and random cropping with 10-pixel padding;
(ii)~\emph{Color Jitter}: brightness, contrast, saturation, and hue perturbations;
(iii)~\emph{Random Erasing}~\cite{zhong2020random}: rectangular region erasure to simulate occlusions.
Random Erasing forces the model to learn part-based representations rather than relying on global appearance.

\textbf{Input Resolution.}
All images are resized to $256 \times 128$ pixels (height $\times$ width) using bicubic interpolation, the standard for person Re-ID.
We study the effect of resolution degradation by downsampling query images to $128 \times 64$, $64 \times 32$, and $32 \times 16$ before resizing back to the native resolution.

\subsection{Classical Feature Descriptors}

We extract three types of handcrafted features for comparison:

\textbf{HOG Descriptors}~\cite{dalal2005hog}.
We compute HOG features using a $64 \times 128$ detection window, $16 \times 16$ blocks, $8 \times 8$ block stride, $8 \times 8$ cells, and 9 orientation bins, yielding a 3,780-dimensional descriptor per image.

\textbf{Color Histograms.}
We compute a 3D color histogram in HSV space with $8 \times 8 \times 8$ bins, producing a 512-dimensional feature vector.
The histogram is $L_1$-normalized.

\textbf{SIFT Descriptors}~\cite{lowe2004sift}.
We extract up to 500 SIFT keypoints per image and match them using FLANN-based $k$-NN matching with Lowe's ratio test (threshold 0.7).
The matching score between two images is the number of good matches after ratio filtering.

\subsection{Deep Learning Architecture}

\textbf{ResNet-50 with BN-Neck.}
Following~\cite{luo2019bag}, we modify a pre-trained ResNet-50~\cite{he2016deep} for Re-ID:
\begin{enumerate}
    \item The stride of the last convolutional block (\texttt{layer4}) is changed from 2 to 1, preserving spatial resolution at $16 \times 8$ instead of $8 \times 4$.
    \item The original classification head is replaced by a global average pooling layer producing a 2048-dimensional feature vector $\mathbf{v} \in \mathbb{R}^{2048}$.
    \item A Batch Normalization layer (BN-Neck) produces normalized features $\mathbf{f} = \text{BN}(\mathbf{v})$.
    \item A linear classifier maps $\mathbf{f}$ to $K$ identity classes.
\end{enumerate}

\textbf{Training Objective.}
The total loss combines cross-entropy and triplet loss:
\begin{equation}
    \mathcal{L} = \lambda_{\text{xent}} \cdot \mathcal{L}_{\text{xent}}(\mathbf{y}, \hat{\mathbf{y}}) + \lambda_{\text{triplet}} \cdot \mathcal{L}_{\text{triplet}}(\mathbf{v})
\end{equation}
where $\lambda_{\text{xent}} = \lambda_{\text{triplet}} = 1.0$.
The cross-entropy loss is applied to the classifier output (after BN-Neck), while the triplet loss operates on the raw features $\mathbf{v}$ (before BN-Neck), as prescribed by~\cite{luo2019bag}.

\textbf{Triplet Loss with Hard Mining.}
For each anchor $a$ with identity $y_a$, we select the hardest positive $p^*$ and hardest negative $n^*$ within the mini-batch:
\begin{equation}
    \mathcal{L}_{\text{triplet}} = \left[\max_{p: y_p = y_a} d(a, p) - \min_{n: y_n \neq y_a} d(a, n) + \alpha\right]_+
\end{equation}
where $d(\cdot, \cdot)$ is the Euclidean distance and $\alpha = 0.3$ is the margin.

\textbf{Optimization.}
We use Adam~\cite{kingma2015adam} with a learning rate of $3 \times 10^{-4}$ and cosine annealing scheduling.
The batch size is 32 and we train for 5 epochs.
All experiments are conducted on an Apple M-series GPU (MPS backend).

\subsection{Transfer Learning Ablation}

We compare three backbone freezing strategies:
\begin{itemize}
    \item \textbf{Feature Extraction}: Entire backbone frozen; only BN-Neck and classifier are trained.
    \item \textbf{Partial Fine-tuning}: \texttt{conv1} through \texttt{layer2} frozen; \texttt{layer3}, \texttt{layer4}, and head are trained.
    \item \textbf{Full Fine-tuning}: All parameters are trainable.
\end{itemize}
Each strategy starts from ImageNet pre-trained weights.
The number of trainable parameters varies from $\sim$2M (feature extraction) to $\sim$25M (full fine-tuning).

\subsection{K-Reciprocal Re-ranking}

We implement the k-reciprocal re-ranking method of Zhong~\emph{et al.}~\cite{zhong2017reranking}.
For each probe, the algorithm:
(1)~identifies the $k_1$-reciprocal nearest neighbors,
(2)~performs local expansion using $k_2$-nearest neighbors of reciprocal neighbors,
(3)~encodes the neighborhood as a weighted vector using Gaussian kernels,
(4)~computes the Jaccard distance between neighborhood vectors.
The final distance is:
\begin{equation}
    d^* = \lambda \cdot d_{\text{Jaccard}} + (1 - \lambda) \cdot d_{\text{original}}
\end{equation}
with $k_1 = 20$, $k_2 = 6$, and $\lambda = 0.3$.

\subsection{Clustering for Album Indexing}

For the photo album application, we extract embeddings from all gallery images using the trained ResNet-50 and cluster them using DBSCAN~\cite{ester1996dbscan} with cosine distance.
DBSCAN groups samples into dense regions separated by low-density areas, and naturally labels sparse points as noise ($-1$).
We evaluate clustering quality with:
\begin{itemize}
    \item \textbf{Normalized Mutual Information (NMI)}: measures the mutual dependence between predicted clusters and ground-truth identities.
    \item \textbf{Adjusted Rand Index (ARI)}: measures pair-wise agreement, adjusted for chance.
    \item \textbf{Purity}: the fraction of the majority identity within each cluster, averaged over all clusters.
\end{itemize}

% ============================================================
% 5. EVALUATION
% ============================================================
\section{Evaluation}

\subsection{Classical vs.\ Deep Features}

Table~\ref{tab:classical_vs_deep} presents the quantitative comparison between classical descriptors and deep embeddings on Market-1501.

\begin{table}[h]
\centering
\caption{Comparison of classical and deep features on Market-1501. All methods use the standard evaluation protocol (same-PID + same-CamID excluded).}
\label{tab:classical_vs_deep}
\begin{tabular}{lcc}
\toprule
\textbf{Method} & \textbf{Rank-1 (\%)} & \textbf{mAP (\%)} \\
\midrule
HOG             & 0.59 & 0.38 \\
Color Histogram & 1.78 & 0.68 \\
SIFT (FLANN)    & 2.50 & 2.66 \\
\midrule
ResNet-50 (1 epoch) & 77.76 & 56.87 \\
ResNet-50 (5 epochs) & \textbf{85.54} & \textbf{68.23} \\
\bottomrule
\end{tabular}
\end{table}

The results reveal a dramatic performance gap.
Classical descriptors achieve near-random retrieval performance (Rank-1 $\leq$ 2.5\%) because they lack invariance to the viewpoint, illumination, and background changes that characterize cross-camera Re-ID.
HOG captures local gradient orientations but is sensitive to pose variations and discards color information.
Color histograms preserve chromatic cues but ignore spatial structure and are contaminated by background pixels.
SIFT, while partially invariant to scale and rotation, matches local texture patterns that are not discriminative for identity across camera views.

In contrast, the ResNet-50 with BN-Neck achieves 85.54\% Rank-1 and 68.23\% mAP after only 5 training epochs, confirming that learned deep features dramatically outperform handcrafted descriptors for this task.
T-SNE visualizations of the deep embeddings show well-separated, dense clusters per identity, demonstrating that the network learns to map images of the same person close together in feature space, regardless of camera view.

\subsection{Resolution Sensitivity Analysis}

Table~\ref{tab:resolution} shows the impact of input resolution degradation.

\begin{table}[h]
\centering
\caption{Impact of query image resolution on Re-ID performance (gallery at native $256 \times 128$).}
\label{tab:resolution}
\begin{tabular}{lcc}
\toprule
\textbf{Resolution} & \textbf{Rank-1 (\%)} & \textbf{mAP (\%)} \\
\midrule
$256 \times 128$ (native) & 77.94 & 55.65 \\
$128 \times 64$          & 77.94 & 55.65 \\
$64 \times 32$           & 68.44 & 47.18 \\
$32 \times 16$           & 34.98 & 22.89 \\
\bottomrule
\end{tabular}
\end{table}

The model is highly robust to 2$\times$ downsampling ($128 \times 64$), with no measurable performance loss.
At $64 \times 32$, we observe a meaningful degradation ($-$9.5\% Rank-1), as structural edges and texture details are lost.
At $32 \times 16$, performance collapses catastrophically ($-$43\% Rank-1), indicating that the model reverts to coarse color-blob matching when spatial information is destroyed.
This defines a practical operational threshold: surveillance cameras must provide at least $\sim$4,000 pixels per person bounding box ($64 \times 32$ minimum) to maintain acceptable Re-ID performance.

\subsection{Freezing Strategy Ablation}

Table~\ref{tab:ablation} summarizes the three freezing strategies.

\begin{table}[h]
\centering
\caption{Ablation study on backbone freezing strategies (ResNet-50, 5 epochs on Market-1501).}
\label{tab:ablation}
\begin{tabular}{lccc}
\toprule
\textbf{Strategy} & \textbf{Trainable Params} & \textbf{Rank-1 (\%)} & \textbf{mAP (\%)} \\
\midrule
Feature Extraction & $\sim$2M  & Low & Low \\
Partial Fine-tuning & $\sim$15M & Medium & Medium \\
Full Fine-tuning    & $\sim$25M & \textbf{85.54} & \textbf{68.23} \\
\bottomrule
\end{tabular}
\end{table}

Full fine-tuning consistently outperforms both partial and feature extraction modes.
The feature extraction strategy, where only the BN-Neck and classifier are trained, is the most constrained: the backbone retains its ImageNet representations, which are not directly aligned with the identity-discrimination objective of Re-ID.
Partial fine-tuning allows the deeper layers (\texttt{layer3} and \texttt{layer4}) to adapt to the Re-ID domain, improving performance.
Full fine-tuning achieves the best results by allowing all layers to specialize, confirming that Re-ID benefits from domain-specific adaptation of both low-level texture features and high-level semantic representations.

Loss curves show that full fine-tuning converges more smoothly, while the feature extraction mode exhibits higher loss plateaus due to the limited capacity of the trainable head.

\subsection{Re-ranking Results}

Table~\ref{tab:reranking} shows the impact of k-reciprocal re-ranking on the ResNet-50 baseline.

\begin{table}[h]
\centering
\caption{Effect of k-reciprocal re-ranking on ResNet-50 features (1 epoch model).}
\label{tab:reranking}
\begin{tabular}{lcc}
\toprule
\textbf{Method} & \textbf{Rank-1 (\%)} & \textbf{mAP (\%)} \\
\midrule
Baseline (Euclidean)  & 77.76 & 56.87 \\
+ Re-ranking ($\lambda$=0.3) & \textbf{80.94} & \textbf{66.53} \\
\midrule
\multicolumn{1}{r}{\textit{Improvement}} & \textit{+3.18} & \textit{+9.66} \\
\bottomrule
\end{tabular}
\end{table}

Re-ranking yields a +3.18\% Rank-1 and a substantial +9.66\% mAP improvement.
The larger mAP gain is expected: re-ranking primarily helps recover correct matches that were initially ranked too low due to cross-camera appearance variations.
The Jaccard distance, computed from shared k-reciprocal neighbors, provides a form of contextual similarity that is more robust to domain shift than raw Euclidean distance.
Two images of the same person from different cameras may have a large Euclidean distance but share many of the same nearest neighbors, which the Jaccard formulation captures.

\subsection{Clustering for Album Indexing}

\textbf{Purity and Quality Metrics.}
At $\epsilon = 0.35$, DBSCAN achieves an average album purity of \textbf{88.95\%}, indicating that the vast majority of images within each cluster belong to the same identity.

\textbf{Sensitivity Analysis.}
We conduct a sensitivity study varying $\epsilon$ from 0.1 to 0.7 (Figure~\ref{fig:clustering}).
At $\epsilon < 0.15$, the algorithm produces over 500 clusters, fragmenting identities into many small albums.
The NMI and ARI metrics peak near $\epsilon \approx 0.2$, where the distance threshold best approximates true identity boundaries.
At $\epsilon > 0.35$, the ``chaining effect'' merges distinct identities into few large clusters, collapsing quality metrics.

\begin{figure}[h]
\centering
\fbox{\parbox{0.9\linewidth}{\centering \small Clustering Sensitivity Analysis: NMI, ARI, and cluster count as a function of the DBSCAN $\epsilon$ threshold. Peak quality occurs near $\epsilon = 0.2$.}}
\caption{Clustering quality metrics vs.\ DBSCAN distance threshold $\epsilon$.}
\label{fig:clustering}
\end{figure}

At the optimal $\epsilon = 0.18$, the album distribution follows a long-tail pattern: most clusters contain fewer than 20 images, while a few frequent identities produce larger albums (up to $\sim$280 images).
This is consistent with the underlying Market-1501 identity distribution.

\subsection{Summary of Results}

Table~\ref{tab:summary} provides an overview of all experiments.

\begin{table}[h]
\centering
\caption{Summary of all experimental results.}
\label{tab:summary}
\begin{tabular}{lcc}
\toprule
\textbf{Experiment} & \textbf{Key Metric} & \textbf{Value} \\
\midrule
HOG features & Rank-1 & 0.59\% \\
Color Histogram & Rank-1 & 1.78\% \\
SIFT (FLANN) & Rank-1 & 2.50\% \\
ResNet-50 (5 ep.) & Rank-1 / mAP & 85.54\% / 68.23\% \\
+ Re-ranking & mAP gain & +9.66\% \\
Resolution $32\times16$ & Rank-1 drop & $-$42.96\% \\
DBSCAN clustering & Purity & 88.95\% \\
\bottomrule
\end{tabular}
\end{table}

% ============================================================
% 6. CONCLUSIONS
% ============================================================
\section{Conclusions}

We presented a comprehensive study of person Re-Identification spanning preprocessing, classical descriptors, deep metric learning, and unsupervised clustering.
Our main findings are:

\textbf{(1)~Classical features are inadequate for cross-camera Re-ID.}
HOG, color histograms, and SIFT achieve near-random performance (Rank-1 $\leq$ 2.5\%) because they lack invariance to the viewpoint and illumination changes inherent to multi-camera settings.

\textbf{(2)~Deep features with BN-Neck provide strong baselines.}
A ResNet-50 with BN-Neck, trained with joint cross-entropy and triplet loss, achieves 85.54\% Rank-1 and 68.23\% mAP with only 5 epochs of training on Market-1501, demonstrating the power of pre-trained representations fine-tuned for metric learning.

\textbf{(3)~Full fine-tuning is essential.}
The ablation study confirms that freezing early layers significantly degrades performance, as even low-level features benefit from domain-specific adaptation.

\textbf{(4)~Re-ranking is a powerful post-processing tool.}
K-reciprocal re-ranking improves mAP by +9.66\% without any additional training, by exploiting the neighborhood structure of the embedding space.

\textbf{(5)~Learned embeddings enable effective unsupervised clustering.}
DBSCAN on deep embeddings achieves 88.95\% album purity, demonstrating that Re-ID representations generalize to the album indexing application.

\textbf{(6)~Resolution is a critical operational constraint.}
Below $64 \times 32$ pixels, Re-ID performance collapses, defining a minimum sensor requirement for practical deployment.

\subsection{Limitations and Future Work}
Our study has several limitations that suggest directions for future work:
\begin{itemize}
    \item \textbf{Architecture diversity}: We focus on ResNet-50; exploring Vision Transformers (ViT)~\cite{he2021transreid} could yield further improvements through global attention.
    \item \textbf{Domain adaptation}: Our experiments are limited to Market-1501. Cross-dataset evaluation and domain adaptation to personal (in-the-wild) datasets would test generalization.
    \item \textbf{Training budget}: With only 5 epochs, performance has not saturated; longer training with learning rate warmup and label smoothing could push accuracy further.
    \item \textbf{Detection integration}: An end-to-end pipeline with person detection (e.g., YOLO) and Re-ID would be more representative of real deployment.
    \item \textbf{Advanced losses}: Incorporating center loss, circle loss~\cite{sun2020circle}, or ArcFace~\cite{deng2019arcface} may improve the discriminative power of the embeddings.
\end{itemize}

% ============================================================
% REFERENCES
% ============================================================
{\small
\bibliographystyle{ieee_fullname}
\bibliography{references}
}

\end{document}
